% Section 1.2 -- Literature Framework
% This fragment is designed to be inserted into the Introduction section
% of the MDPI manuscript with:
% % Section 1.2 -- Literature Framework
% This fragment is designed to be inserted into the Introduction section
% of the MDPI manuscript with:
% % Section 1.2 -- Literature Framework
% This fragment is designed to be inserted into the Introduction section
% of the MDPI manuscript with:
% % Section 1.2 -- Literature Framework
% This fragment is designed to be inserted into the Introduction section
% of the MDPI manuscript with:
% \input{Section_1_2_Literature_Framework.tex}

\subsection{Literature Framework}
\label{sec:literature-framework}

The central question in key-sector research is whether the sector with the greatest number of economic linkages is also the sector that yields the highest social return on public resources. The classical forward--backward linkage approach largely conflates these two attributes \citep{ref-hirschman1958,ref-rasmussen1956,ref-chenery-watanabe1958}. This equivalence is nevertheless contested because different key-sector measures generate different rankings \citep{ref-temurshoev-oosterhaven2014}, hypothetical extraction and direct linkages do not produce the same results \citep{ref-pereira-lopez2021,ref-bhusal2025}, and the selected sector changes with the objective---production, the environment, regional development, or foreign trade \citep{ref-alcantara-padilla2020,ref-ernawati2024,ref-ojaleye-narayanan2022,ref-nugroho2021}. Comparative network analysis that includes Turkey \citep{ref-jahangard-keshtvarz2012} and a recent input--output network application for Turkey \citep{ref-cirpici2024} suggest that linkages provide indispensable starting information, but not an automatic, policy-objective-independent list of sectors to support.

This debate has become more concrete with the production-networks approach. Hub/authority, eigenvector, PageRank, betweenness, and random-walk measures consider not only the volume of sectors' purchases and sales, but also the nodes through which shocks are transmitted and the intermediate positions in which they become concentrated \citep{ref-alatriste2015,ref-alatriste-lugo2022,ref-depaolis2022,ref-an2024}. Studies of the EU, Ecuador, Korea, the world input--output network, and multi-country sensitivity show that central sectors often overlap only partially with those selected by traditional linkage measures \citep{ref-ramirez-alvarez2023,ref-lee-ishiro2025,ref-tran2017,ref-wirkierman2021}. This evidence does not diminish the value of centrality; rather, it positions centrality as an indicator of propagation, vulnerability, and structural intermediation. Dynamic-network and disaggregated-shock studies likewise show that the level of sectoral detail and the linkage structure affect the results \citep{ref-guerrieri2014,ref-roson-sartori2016,ref-jiang2024,ref-kharazi2025}.

The general-equilibrium literature makes clear why centrality alone is insufficient for policy. Under efficient conditions, Domar weights and input--output multipliers may summarize first-order aggregate effects; once substitution, distortions, complementarities, and factor mobility enter the analysis, however, the impact of sectoral shocks changes nonlinearly \citep{ref-fadinger2022,ref-baqaee-farhi2019b,ref-baqaee-farhi2020a,ref-baqaee-farhi2020b,ref-bahal-lenzo2023}. Applied CGE studies make visible the growth, welfare, price, and structural-transformation channels of productivity or technology shocks \citep{ref-cutler-davies2010,ref-burnett2012,ref-antoszewski2019,ref-khan2020,ref-blackburn-moreno2021,ref-jin-li2025}, while particularly emphasizing the sensitivity of CGE--sectoral model linkages to their assumptions \citep{ref-delzeit2020,ref-bolarinwa2024,ref-odior-arinze2022}. Yet most of these studies do not directly test whether network centrality predicts size-adjusted, multichannel CGE effects out of sample within the same universe of sectors.

The industrial-policy literature turns this gap into a normative warning: support priorities should derive not only from topological position, but also from where distortions accumulate, the cost of intervention, and the targeted outcome \citep{ref-liu2019,ref-martinez-orrillo2024}. Multi-criteria ranking tools are therefore useful, but they cannot substitute for a decision rule. Entropy weighting, PCA, TOPSIS, fuzzy multi-criteria decision making, and machine-learning applications make it possible to combine different indicators transparently \citep{ref-andre-cardenete2009,ref-jin2022,ref-tsoulfidis-athanasiadis2022,ref-calik2018,ref-turegun2022,ref-jamwal2021,ref-ionascu2024,ref-hlushko2025,ref-faridzad2025}. It must nevertheless also be shown which indicator is selected for which objective and whether the ranking persists across alternative reasonable specifications.

This study brings these two strands of literature together around the same empirical question for Turkey. Within the same 62-sector data universe, it compares four structural dimensions of the network with a size-controlled and channel-balanced EIPI derived from the SAM/CGE results of 62 separate TFP shocks. It evaluates their relationship not only through correlations, but also through classification agreement, LOOCV prediction, a Pareto--balance filter, and robustness across 96 specifications. It thereby tests which types of outcomes contain information from centrality, when that information is confounded by sector size, and to what extent it diverges from explicit policy preferences. This design treats network mapping, CGE shock analysis, and composite rankings---often conducted separately in the literature---not as interchangeable validations, but as complementary bodies of evidence that test one another.



\subsection{Literature Framework}
\label{sec:literature-framework}

The central question in key-sector research is whether the sector with the greatest number of economic linkages is also the sector that yields the highest social return on public resources. The classical forward--backward linkage approach largely conflates these two attributes \citep{ref-hirschman1958,ref-rasmussen1956,ref-chenery-watanabe1958}. This equivalence is nevertheless contested because different key-sector measures generate different rankings \citep{ref-temurshoev-oosterhaven2014}, hypothetical extraction and direct linkages do not produce the same results \citep{ref-pereira-lopez2021,ref-bhusal2025}, and the selected sector changes with the objective---production, the environment, regional development, or foreign trade \citep{ref-alcantara-padilla2020,ref-ernawati2024,ref-ojaleye-narayanan2022,ref-nugroho2021}. Comparative network analysis that includes Turkey \citep{ref-jahangard-keshtvarz2012} and a recent input--output network application for Turkey \citep{ref-cirpici2024} suggest that linkages provide indispensable starting information, but not an automatic, policy-objective-independent list of sectors to support.

This debate has become more concrete with the production-networks approach. Hub/authority, eigenvector, PageRank, betweenness, and random-walk measures consider not only the volume of sectors' purchases and sales, but also the nodes through which shocks are transmitted and the intermediate positions in which they become concentrated \citep{ref-alatriste2015,ref-alatriste-lugo2022,ref-depaolis2022,ref-an2024}. Studies of the EU, Ecuador, Korea, the world input--output network, and multi-country sensitivity show that central sectors often overlap only partially with those selected by traditional linkage measures \citep{ref-ramirez-alvarez2023,ref-lee-ishiro2025,ref-tran2017,ref-wirkierman2021}. This evidence does not diminish the value of centrality; rather, it positions centrality as an indicator of propagation, vulnerability, and structural intermediation. Dynamic-network and disaggregated-shock studies likewise show that the level of sectoral detail and the linkage structure affect the results \citep{ref-guerrieri2014,ref-roson-sartori2016,ref-jiang2024,ref-kharazi2025}.

The general-equilibrium literature makes clear why centrality alone is insufficient for policy. Under efficient conditions, Domar weights and input--output multipliers may summarize first-order aggregate effects; once substitution, distortions, complementarities, and factor mobility enter the analysis, however, the impact of sectoral shocks changes nonlinearly \citep{ref-fadinger2022,ref-baqaee-farhi2019b,ref-baqaee-farhi2020a,ref-baqaee-farhi2020b,ref-bahal-lenzo2023}. Applied CGE studies make visible the growth, welfare, price, and structural-transformation channels of productivity or technology shocks \citep{ref-cutler-davies2010,ref-burnett2012,ref-antoszewski2019,ref-khan2020,ref-blackburn-moreno2021,ref-jin-li2025}, while particularly emphasizing the sensitivity of CGE--sectoral model linkages to their assumptions \citep{ref-delzeit2020,ref-bolarinwa2024,ref-odior-arinze2022}. Yet most of these studies do not directly test whether network centrality predicts size-adjusted, multichannel CGE effects out of sample within the same universe of sectors.

The industrial-policy literature turns this gap into a normative warning: support priorities should derive not only from topological position, but also from where distortions accumulate, the cost of intervention, and the targeted outcome \citep{ref-liu2019,ref-martinez-orrillo2024}. Multi-criteria ranking tools are therefore useful, but they cannot substitute for a decision rule. Entropy weighting, PCA, TOPSIS, fuzzy multi-criteria decision making, and machine-learning applications make it possible to combine different indicators transparently \citep{ref-andre-cardenete2009,ref-jin2022,ref-tsoulfidis-athanasiadis2022,ref-calik2018,ref-turegun2022,ref-jamwal2021,ref-ionascu2024,ref-hlushko2025,ref-faridzad2025}. It must nevertheless also be shown which indicator is selected for which objective and whether the ranking persists across alternative reasonable specifications.

This study brings these two strands of literature together around the same empirical question for Turkey. Within the same 62-sector data universe, it compares four structural dimensions of the network with a size-controlled and channel-balanced EIPI derived from the SAM/CGE results of 62 separate TFP shocks. It evaluates their relationship not only through correlations, but also through classification agreement, LOOCV prediction, a Pareto--balance filter, and robustness across 96 specifications. It thereby tests which types of outcomes contain information from centrality, when that information is confounded by sector size, and to what extent it diverges from explicit policy preferences. This design treats network mapping, CGE shock analysis, and composite rankings---often conducted separately in the literature---not as interchangeable validations, but as complementary bodies of evidence that test one another.



\subsection{Literature Framework}
\label{sec:literature-framework}

The central question in key-sector research is whether the sector with the greatest number of economic linkages is also the sector that yields the highest social return on public resources. The classical forward--backward linkage approach largely conflates these two attributes \citep{ref-hirschman1958,ref-rasmussen1956,ref-chenery-watanabe1958}. This equivalence is nevertheless contested because different key-sector measures generate different rankings \citep{ref-temurshoev-oosterhaven2014}, hypothetical extraction and direct linkages do not produce the same results \citep{ref-pereira-lopez2021,ref-bhusal2025}, and the selected sector changes with the objective---production, the environment, regional development, or foreign trade \citep{ref-alcantara-padilla2020,ref-ernawati2024,ref-ojaleye-narayanan2022,ref-nugroho2021}. Comparative network analysis that includes Turkey \citep{ref-jahangard-keshtvarz2012} and a recent input--output network application for Turkey \citep{ref-cirpici2024} suggest that linkages provide indispensable starting information, but not an automatic, policy-objective-independent list of sectors to support.

This debate has become more concrete with the production-networks approach. Hub/authority, eigenvector, PageRank, betweenness, and random-walk measures consider not only the volume of sectors' purchases and sales, but also the nodes through which shocks are transmitted and the intermediate positions in which they become concentrated \citep{ref-alatriste2015,ref-alatriste-lugo2022,ref-depaolis2022,ref-an2024}. Studies of the EU, Ecuador, Korea, the world input--output network, and multi-country sensitivity show that central sectors often overlap only partially with those selected by traditional linkage measures \citep{ref-ramirez-alvarez2023,ref-lee-ishiro2025,ref-tran2017,ref-wirkierman2021}. This evidence does not diminish the value of centrality; rather, it positions centrality as an indicator of propagation, vulnerability, and structural intermediation. Dynamic-network and disaggregated-shock studies likewise show that the level of sectoral detail and the linkage structure affect the results \citep{ref-guerrieri2014,ref-roson-sartori2016,ref-jiang2024,ref-kharazi2025}.

The general-equilibrium literature makes clear why centrality alone is insufficient for policy. Under efficient conditions, Domar weights and input--output multipliers may summarize first-order aggregate effects; once substitution, distortions, complementarities, and factor mobility enter the analysis, however, the impact of sectoral shocks changes nonlinearly \citep{ref-fadinger2022,ref-baqaee-farhi2019b,ref-baqaee-farhi2020a,ref-baqaee-farhi2020b,ref-bahal-lenzo2023}. Applied CGE studies make visible the growth, welfare, price, and structural-transformation channels of productivity or technology shocks \citep{ref-cutler-davies2010,ref-burnett2012,ref-antoszewski2019,ref-khan2020,ref-blackburn-moreno2021,ref-jin-li2025}, while particularly emphasizing the sensitivity of CGE--sectoral model linkages to their assumptions \citep{ref-delzeit2020,ref-bolarinwa2024,ref-odior-arinze2022}. Yet most of these studies do not directly test whether network centrality predicts size-adjusted, multichannel CGE effects out of sample within the same universe of sectors.

The industrial-policy literature turns this gap into a normative warning: support priorities should derive not only from topological position, but also from where distortions accumulate, the cost of intervention, and the targeted outcome \citep{ref-liu2019,ref-martinez-orrillo2024}. Multi-criteria ranking tools are therefore useful, but they cannot substitute for a decision rule. Entropy weighting, PCA, TOPSIS, fuzzy multi-criteria decision making, and machine-learning applications make it possible to combine different indicators transparently \citep{ref-andre-cardenete2009,ref-jin2022,ref-tsoulfidis-athanasiadis2022,ref-calik2018,ref-turegun2022,ref-jamwal2021,ref-ionascu2024,ref-hlushko2025,ref-faridzad2025}. It must nevertheless also be shown which indicator is selected for which objective and whether the ranking persists across alternative reasonable specifications.

This study brings these two strands of literature together around the same empirical question for Turkey. Within the same 62-sector data universe, it compares four structural dimensions of the network with a size-controlled and channel-balanced EIPI derived from the SAM/CGE results of 62 separate TFP shocks. It evaluates their relationship not only through correlations, but also through classification agreement, LOOCV prediction, a Pareto--balance filter, and robustness across 96 specifications. It thereby tests which types of outcomes contain information from centrality, when that information is confounded by sector size, and to what extent it diverges from explicit policy preferences. This design treats network mapping, CGE shock analysis, and composite rankings---often conducted separately in the literature---not as interchangeable validations, but as complementary bodies of evidence that test one another.



\subsection{Literature Framework}
\label{sec:literature-framework}

The central question in key-sector research is whether the sector with the greatest number of economic linkages is also the sector that yields the highest social return on public resources. The classical forward--backward linkage approach largely conflates these two attributes \citep{ref-hirschman1958,ref-rasmussen1956,ref-chenery-watanabe1958}. This equivalence is nevertheless contested because different key-sector measures generate different rankings \citep{ref-temurshoev-oosterhaven2014}, hypothetical extraction and direct linkages do not produce the same results \citep{ref-pereira-lopez2021,ref-bhusal2025}, and the selected sector changes with the objective---production, the environment, regional development, or foreign trade \citep{ref-alcantara-padilla2020,ref-ernawati2024,ref-ojaleye-narayanan2022,ref-nugroho2021}. Comparative network analysis that includes Turkey \citep{ref-jahangard-keshtvarz2012} and a recent input--output network application for Turkey \citep{ref-cirpici2024} suggest that linkages provide indispensable starting information, but not an automatic, policy-objective-independent list of sectors to support.

This debate has become more concrete with the production-networks approach. Hub/authority, eigenvector, PageRank, betweenness, and random-walk measures consider not only the volume of sectors' purchases and sales, but also the nodes through which shocks are transmitted and the intermediate positions in which they become concentrated \citep{ref-alatriste2015,ref-alatriste-lugo2022,ref-depaolis2022,ref-an2024}. Studies of the EU, Ecuador, Korea, the world input--output network, and multi-country sensitivity show that central sectors often overlap only partially with those selected by traditional linkage measures \citep{ref-ramirez-alvarez2023,ref-lee-ishiro2025,ref-tran2017,ref-wirkierman2021}. This evidence does not diminish the value of centrality; rather, it positions centrality as an indicator of propagation, vulnerability, and structural intermediation. Dynamic-network and disaggregated-shock studies likewise show that the level of sectoral detail and the linkage structure affect the results \citep{ref-guerrieri2014,ref-roson-sartori2016,ref-jiang2024,ref-kharazi2025}.

The general-equilibrium literature makes clear why centrality alone is insufficient for policy. Under efficient conditions, Domar weights and input--output multipliers may summarize first-order aggregate effects; once substitution, distortions, complementarities, and factor mobility enter the analysis, however, the impact of sectoral shocks changes nonlinearly \citep{ref-fadinger2022,ref-baqaee-farhi2019b,ref-baqaee-farhi2020a,ref-baqaee-farhi2020b,ref-bahal-lenzo2023}. Applied CGE studies make visible the growth, welfare, price, and structural-transformation channels of productivity or technology shocks \citep{ref-cutler-davies2010,ref-burnett2012,ref-antoszewski2019,ref-khan2020,ref-blackburn-moreno2021,ref-jin-li2025}, while particularly emphasizing the sensitivity of CGE--sectoral model linkages to their assumptions \citep{ref-delzeit2020,ref-bolarinwa2024,ref-odior-arinze2022}. Yet most of these studies do not directly test whether network centrality predicts size-adjusted, multichannel CGE effects out of sample within the same universe of sectors.

The industrial-policy literature turns this gap into a normative warning: support priorities should derive not only from topological position, but also from where distortions accumulate, the cost of intervention, and the targeted outcome \citep{ref-liu2019,ref-martinez-orrillo2024}. Multi-criteria ranking tools are therefore useful, but they cannot substitute for a decision rule. Entropy weighting, PCA, TOPSIS, fuzzy multi-criteria decision making, and machine-learning applications make it possible to combine different indicators transparently \citep{ref-andre-cardenete2009,ref-jin2022,ref-tsoulfidis-athanasiadis2022,ref-calik2018,ref-turegun2022,ref-jamwal2021,ref-ionascu2024,ref-hlushko2025,ref-faridzad2025}. It must nevertheless also be shown which indicator is selected for which objective and whether the ranking persists across alternative reasonable specifications.

This study brings these two strands of literature together around the same empirical question for Turkey. Within the same 62-sector data universe, it compares four structural dimensions of the network with a size-controlled and channel-balanced EIPI derived from the SAM/CGE results of 62 separate TFP shocks. It evaluates their relationship not only through correlations, but also through classification agreement, LOOCV prediction, a Pareto--balance filter, and robustness across 96 specifications. It thereby tests which types of outcomes contain information from centrality, when that information is confounded by sector size, and to what extent it diverges from explicit policy preferences. This design treats network mapping, CGE shock analysis, and composite rankings---often conducted separately in the literature---not as interchangeable validations, but as complementary bodies of evidence that test one another.

